\begin{examplebox}
\textbf{例1：从椭圆方程求焦点。} $x^2/25+y^2/16=1$ 中 $a=5,b=4$，
$c=\sqrt{25-16}=3$，焦点为 $\boxed{(\pm3,0)}$。
回验：短轴端点 $(0,4)$ 到两焦点距离均为 $5$，和为 $10$。

\medskip
\textbf{例2：双曲线不按分母大小读 $a$。}
$x^2/9-y^2/16=1$ 的正项给出 $a=3$，负项给出 $b=4$，故
$c=5$，实轴长 $6$、虚轴长 $8$、焦点 $\boxed{(\pm5,0)}$。
回验：右顶点 $(3,0)$ 到焦点距离 $8,2$，差为 $6=2a$。

\medskip
\textbf{例3：由轴长、焦距反求方程，中心条件不可少。}
\textbf{中心在原点}、焦点在 $x$ 轴的椭圆，焦距 $8$、长轴长 $10$。
先取 $c=4,a=5$，再算 $b^2=25-16=9$，得 $\boxed{x^2/25+y^2/9=1}$。
回验：$a>c$，长轴 $2a=10$，焦距 $2c=8$。
\begin{center}\begin{tikzpicture}[font=\small,>=stealth]
\draw[black!30,rounded corners=2pt] (0,0) rectangle (14.4,8.2);
\node[anchor=west,font=\small\bfseries] at (.5,7.72) {焦点都在横轴上，中心也可能不在原点};

\begin{scope}[shift={(6.4,4.2)},scale=.53]
\draw[black!35,->] (-6,0)--(8,0) node[right] {$x$};
\draw[blue,very thick] (0,0) ellipse (5 and 3);
\draw[orange!85!black,very thick,dashed] (2,0) ellipse (5 and 3);
\fill[blue] (-4,0) circle (3pt);\fill[blue] (4,0) circle (3pt);
\fill[orange!85!black] (-2,0) circle (3pt);\fill[orange!85!black] (6,0) circle (3pt);
\draw[black!50,->] (0,4.1)--(2,4.1) node[midway,above=5pt] {平移 $2$};
\end{scope}
\node[blue] at (7.2,1.15) {$x^2/25+y^2/9=1$；中心 $(0,0)$};
\node[orange!85!black] at (7.2,1.8) {$(x-2)^2/25+y^2/9=1$；中心 $(2,0)$};

\node[anchor=west] at (.7,.55) {同尺度；两图长轴长 $10$、焦距 $8$ 均相同，故写原点标准式必须说明中心};
\end{tikzpicture}\end{center}

若只说“焦点在 $x$ 轴上”，中心还可沿横轴平移；图中另一条曲线也是同样的轴长、焦距，方程却不同。

\medskip
\textbf{边界反例。} 焦点 $(\pm3,0)$ 的距离差不可能是 $8$，因为
$|PF_1-PF_2|\le F_1F_2=6$。不能硬取负的 $b^2$ 再改成正数写方程。
\end{examplebox}
\clearpage
